\section{Background and Challenges}
\subsection{Neural Language Model Pretraining} 

Pretrained language models have become an indispensable part of NLP researchers' toolkits. Leveraging large corpus pretraining to learn robust neural representations of language is an active area of research that has spanned the past decade. Early examples of pretraining and transferring neural representations of language demonstrated that pretrained word embedding tables improve downstream task results compared to word embedding tables learned from scratch~\cite{Mikolov2013,Pennington2014,Turian2010}. 
Later work advanced research in this area by learning and transferring neural models that capture contextual representations of words~\citep{context2vec, Cove, ELMo, Radford2017Sentiment, Radford2019GPT2}. Recent parallel work \citep{Le2016seq2seqtransfer, Howard2018ULMFIT, Radford2018GPT, devlin2018bert, roberta, transformerxl,xlnet,mtdnn,ALBERT2019} further builds upon these ideas by not just transferring the language model to extract contextual word representations, but by also finetuning the language model in an end to end fashion on downstream tasks. Through these works, the state of the art has advanced from transferring just word embedding tables to transferring entire multi-billion parameter language models. This progression of methods has necessitated the need for hardware, systems techniques, and frameworks that are able to operate efficiently at scale and satisfy increasing computational needs. Our work aims to provide the tools necessary to take another step forward in this trend.

\subsection{Transformer Language Models and Multi-Head Attention} 

Current work in NLP trends towards using {\it transformer} models~\cite{Transformer} due to their superior accuracy and compute efficiency. The original transformer formulation was designed as a machine translation architecture that transforms an input sequence into another output sequence using two parts, an {\it Encoder} and {\it Decoder}. However, recent work leveraging transformers for language modeling such as BERT \cite{devlin2018bert} and GPT-2 \cite{Radford2019GPT2} use only the {\it Encoder} or {\it Decoder} depending on their needs. This work explores both a decoder architecture, GPT-2, and an encoder architecture, BERT. 

\begin{figure}[t!]
\vspace{-2mm}
\begin{center}
  \includegraphics[scale=0.24]{./transformer-general.jpeg}
  \vspace{-2mm}
  \caption{Transformer Architecture. Purple blocks correspond to fully connected layers. Each blue block represents a single transformer layer that is replicated N times.}
  \label{fig:gpt2-transformer}
\end{center}
\vspace{-6mm}
\end{figure}

Figure \ref{fig:gpt2-transformer} shows a schematic diagram of the model we used. We refer the reader to prior work for a detailed description of the model architecture~\citep{Transformer,devlin2018bert,Radford2019GPT2}. It is worthwhile to mention that both GPT-2 and BERT use GeLU \citep{gelu} nonlinearities and layer normalization \citep{layernorm} to the input of the multi-head attention and feed forward layers, whereas the original transformer \citep{Transformer} uses ReLU nonlinearities and applies layer normalization to outputs.

\subsection{Data and Model Parallelism in Deep Learning} 

There are two central paradigms for scaling out deep neural network training to numerous hardware accelerators: data parallelism \cite{DataParallel1990} where a training minibatch is split across multiple workers, and model parallelism in which the memory usage and computation of a model is distributed across multiple workers. By increasing the minibatch size proportionally to the number of available workers (i.e. \emph{weak scaling}), one observes near linear scaling in training data throughput. However, large batch training introduces complications into the optimization process that can result in reduced accuracy or longer time to convergence, offsetting the benefit of increased training throughput \cite{LargeBatch2016}. Further research \cite{Goyal2017, Lars2017, Lamb2019} has developed techniques to mitigate these effects and drive down the training time of large neural networks. To scale out training even further, parallel work \cite{activation_checkpointing} has combined data parallelism with activation checkpointing: recomputing activations in the backward pass without storing them in the forward pass to reduce memory requirements. 

However, these techniques have one fundamental limitation in the problem size they can tackle: the model must fit entirely on one worker. With language models of increasing size and complexity like BERT and GPT-2, neural networks have approached the memory capacity of modern hardware accelerators. One solution to this problem is to employ parameter sharing to reduce the memory footprint of the model \cite{ALBERT2019}, but this limits the overall capacity of the model. Our approach is to utilize model parallelism to split the model across multiple accelerators. This not only alleviates the memory pressure, but also increases the amount of parallelism independently of the microbatch size.

Within model parallelism, there are two further paradigms: layer-wise pipeline parallelism, and more general distributed tensor computation. In pipeline model parallelism, groups of operations are performed on one device before the outputs are passed to the next device in the pipeline where a different group of operations are performed. Some approaches \cite{PipeDream2018, DualPipe2018} use a parameter server \cite{ParamServer2014} in conjunction with pipeline parallelism. However these suffer from inconsistency issues. The GPipe framework for TensorFlow \cite{GPipe} overcomes this inconsistency issue by using synchronous gradient decent. This approach requires additional logic to handle the efficient pipelining of these communication and computation operations, and suffers from pipeline bubbles that reduce efficiency, or changes to the optimizer itself which impact accuracy.

Distributed tensor computation is an orthogonal and more general approach that partitions a tensor operation across multiple devices to accelerate computation or increase model size. FlexFlow \cite{FlexFlow2018}, a deep learning framework orchestrating such parallel computation, provides a method to pick the best parallelization strategy. Recently, Mesh-TensorFlow \cite{mesh_tf} introduced a language for specifying a general class of distributed tensor computations in TensorFlow \cite{tensorflow2015-whitepaper}. The parallel dimensions are specified in the language by the end user and the resulting graph is compiled with proper collective primitives. We utilize similar insights to those leveraged in Mesh-TensorFlow and exploit parallelism in computing the transformer's attention heads to parallelize our transformer model. However, rather than implementing a framework and compiler for model parallelism, we make only a few targeted modifications to existing PyTorch transformer implementations. Our approach is simple, does not require any new compiler or code re-writing, and can be fully implemented by inserting a few simple primitives, as described in the next section.